% Shared preamble for the Week 1 student pages.
\usepackage[letterpaper,left=0.55in,right=0.55in,top=0.85in,bottom=0.8in,
            headheight=18pt,headsep=14pt,footskip=28pt]{geometry}
\usepackage[T1]{fontenc}
\usepackage[default]{andika}
\usepackage{tikz}
\usetikzlibrary{arrows.meta}
\usepackage{fancyhdr}
\usepackage{xcolor}

\definecolor{gridgray}{gray}{0.60}
\definecolor{chainfill}{gray}{0.78}
\definecolor{holefill}{gray}{0.08}
\definecolor{blockG}{HTML}{3DAE4B}
\definecolor{blockB}{HTML}{3F6FD8}
\definecolor{blockR}{HTML}{E2463B}
\definecolor{blockY}{HTML}{F6D12E}
\definecolor{blockP}{HTML}{8A4FC0}
\colorlet{fillG}{blockG!25}
\colorlet{fillB}{blockB!18}
\colorlet{fillR}{blockR!25}
\colorlet{fillY}{blockY!35}
\colorlet{fillP}{blockP!25}
\tikzset{writeline/.style={black!45,line width=0.5pt}}

\pagestyle{fancy}
\fancyhf{}
\renewcommand{\headrulewidth}{0.4pt}
\renewcommand{\footrulewidth}{0pt}
\lfoot{Bellingham Math Circle / Week 1}
\rfoot{\thepage}

\setlength{\parindent}{0pt}
\setlength{\parskip}{0pt}
\raggedright

% "Problem N:" label followed by ordinary sentences.
\newcommand{\problem}[1]{\textbf{Problem #1:}\ }
\newcommand{\fig}[1]{\par\noindent\input{figs/#1}\par}

% A block picture followed by its name, for the "which blocks" line of a problem.
\newcommand{\blk}[2]{\mbox{\input{figs/icon_#1}\ \ #2}}
\newcommand{\blkgap}{\hspace{0.45in}}
